\section{引言}

小型无人机已广泛应用于物流、巡检、应急响应和低空感知，同时也使重点空域产生了稳定、低成本的视觉监测需求。与通用目标检测不同，远距离无人机在缩放后往往不足 $16\times16$ 像素，快速运动、失焦和压缩削弱轮廓，云层、鸟类、建筑和植被又会形成相似响应。检测器既要提高召回率、避免漏警，又要抑制背景误检，并满足边缘平台的存储和延迟约束。现有综述将这种精度与效率矛盾视为视觉反无人机系统的核心问题之一\cite{yasmeen2025review,wang2021smallobjectsurvey}。

YOLO 与 SSD 奠定了单阶段密集预测的实时检测范式\cite{redmon2016yolo,liu2016ssd}，FPN 和 PANet 进一步增强高层语义与浅层空间信息的交换\cite{lin2017fpn,liu2018panet}。FCOS、Focal Loss 与 GFL 分别从无锚预测、类别不平衡和质量感知置信度角度提升密集检测\cite{tian2019fcos,lin2017focal,li2020gfl}。然而，直接增加高分辨率分支或叠加注意力模块会提高显存访问与算子调度开销。EfficientDet、YOLOv4 和 YOLOv7 的实践表明，实际速度取决于骨干、颈部、训练策略和部署图的协同设计，而非仅由 FLOPs 决定\cite{tan2020efficientdet,bochkovskiy2020yolov4,wang2023yolov7}。

本文的出发点是分类与定位并不需要完全相同的特征子空间。分类偏好语义不变性和类间判别性，定位则必须保留对边界和位移敏感的信息。TSD 表明共享候选与表征可能造成任务冲突\cite{wu2020tsd}，TOOD 则说明任务交互需要与任务专属特征及对齐机制并存\cite{feng2021tood}。完整解耦头能够减少直接特征竞争，却会复制代价较高的空间计算；完全共享头更为紧凑，但两项任务受到同一表征预算约束。因此，本文采用受约束的折中方案：保留两任务共用的紧凑基底，仅为分类分配少量附加表征。

基于上述分析，本文提出 Shared-Representation and Private-Rank Allocation YOLO（SRPA-YOLO），作为 YOLOv8 单阶段框架的结构改进，并以 YOLOv8n 为归因精度、参数量与延迟变化的主要基线。模型通过轻量双向颈部保留 P2--P4 特征，每一尺度的检测头包含 64 通道共享 RepConv 和 16 通道分类私有 RepConv。私有输出投影初始化为零，从而保持共享检测器初始的输入--输出函数。部署前采用 RepVGG 的结构重参数化思想\cite{ding2021repvgg}，将两条等效 stem 卷积核沿输出通道拼接，并把任务投影合并为单路径。模型仅在训练期扩展分类容量，推理时不保留并行私有分支。

训练采用两个阶段。阶段 I 利用标准监督、分类响应蒸馏和离散框分布蒸馏学习共享检测器\cite{zheng2022ld}，其中置信度加权项强调可靠的教师定位响应；阶段 II 从函数保持的零残差状态出发，仅优化分类私有 stem 与投影。在该约束下，既有边界框回归函数保持不变，分类获得附加表征容量。本文的方法贡献在于任务专属表征分配、函数保持初始化与精确部署融合的组合；任务解耦、零初始化、蒸馏和结构重参数化本身均已有研究基础。具体贡献如下：
\begin{enumerate}
\item 将预测头容量表述为共享表征与分类私有表征的分配问题。64+16 结构使边界框输出在计算图上直接独立于私有表征，同时允许分类使用共享与私有表征的组合行空间。
\item 构造函数保持的两阶段实现：共享检测器收敛后，从零输出私有投影开始细化，仅增加分类容量，而阶段 II 的边界框回归映射保持固定。
\item 给出精确部署融合公式，将训练期共享/私有路径转换为每尺度一个 $3\times3$ stem 和一个 $1\times1$ 输出层，在浮点舍入误差范围内保持输出，并消除私有分支推理开销。
\item 采用验证集选模、尺度与 IoU 分解、阶段 II 的 40/41/42 匹配随机种子以及同进程配对延迟进行评估。相对 YOLOv8n，最终 checkpoint 的 $P$、$R$、AP@0.50、\mapall{} 和 AP@0.75 分别提高 3.088、8.242、5.057、6.895 和 9.750 个百分点，部署参数量减少 28.0\%。
\end{enumerate}
